% vim: set foldmethod=marker foldlevel=0:

\documentclass[a4paper]{article}
\usepackage[UKenglish]{babel}

\usepackage{preamble}

\fancyhead[L]{MA263 Assignment 3}
\title{MA263 Multivariable Analysis, Assignment 3}
\colorlet{questionbodycolor}{cyan!50}

\begin{document}

\maketitle

\setlength{\parindent}{0em}
\setlength{\parskip}{1em}

% {{{ Q1
\question{1}

\begin{questionbody}
Let $M \subset \R^k$ and $N \subset \R^\ell$ be smooth manifolds of dimensions $m$ and $n$ respectively.

We say that a function $f : M \to N$ is smooth if it is locally the restriction to $M$ of a smooth function. That is, for every $a \in M$, there exists open $U \subset \R^k$ containing $a$ and a smooth function $g : U \to \R^\ell$ with $g(x) = f(x)$ for every $x \in M \cap U$.

\begin{enumerate}[(a)]
\item Define, for $a \in M$, the derivative $D f(a) : T_a M \to T_{f(a)} N$ by $D f(a) v = D g(a) v$ for $v \in T_a M$, where $g$ is as above.

Show that this is well-define, meaning it is independent of $g$ and maps into $T_{f(a)} N$.

\item We say that $f$ is a diffeomorphism if it is invertible and its inverse is also smooth. In this case we say that $M$ and $N$ are diffeomorphic.

Prove that diffeomorphic smooth manifolds must have the same dimension.
\end{enumerate}
\end{questionbody}

\subsection{~} % 1.a

Answer

\subsection{~} % 1.b

Answer

% }}}

% {{{ Q2
\newquestion{2}

\begin{questionbody}
Let $U \subset \R^m$ and let $r : U \to \R^n$ be a parametrisation of a manifold $M$.

\begin{enumerate}[(a)]
\item Show that the vectors $\partial_i r(x)$ form a basis of the tangent space $T_{r(x)} M$.

\item Let $m = n - 1$ and consider the Jacobian $\partial r(x)$, which is an $n \times (n-1)$ matrix with the vectors $\partial_i r(x)$ as its columns.

Define a vector $v = (v_1, \dotsc, v_n) \in \R^n$ by $v_i = {(-1)}^{i+n} M_i$, where $M_i$ is the determinant of the matrix obtained by deleting row $i$ from the Jacobian.

Show that this is a non-zero vector normal to $M$ at $r(x)$.

\textit{Hint}: Using cofactor expansion, you can realise $\angb{v, w}$ as the determinant of the matrix obtained from adjoining $\partial r(x)$ to $w$. Once you have justified this observation, all that's left is some linear algebra.
\end{enumerate}
\end{questionbody}

\subsection{~} % 2.a

Answer

\subsection{~} % 2.b

Answer

% }}}

% {{{ Q3
\newquestion{3}

\begin{questionbody}
Consider the function \[
f(x, y, z) = 3x^2 + 2xy + 2y^2 + z^2
\] restricted to the subset of $\R^3$ determined by \[
2x + 4y = 5, \quad x - y + 2z = 3.
\]

\begin{enumerate}[(a)]
\item Argue that this function indeed attains a minimum on this domain.

\textit{Hint}: First show that $f(x, y, z) \ge \|(x, y, z)\|_2^2$.

\item Find the point at which this minimum is attained, justifying the use of any theorems you apply.

\textit{Note}: During calculations, you will likely get a system of linear equations to solve. State clearly what this system is, but you don't need to submit the calculations for solving the system---just the solution.
\end{enumerate}
\end{questionbody}

\subsection{~} % 3.a

Answer

\subsection{~} % 3.b

Answer

% }}}

\end{document}
